% ==============================================================
%  Chapter 8: AI as a Problem-Solving Partner
% ==============================================================
\chapter{AI as a Problem-Solving Partner}
\label{ch:problem-solving}

Problem-solving is not a single skill --- it is a set of habits: decomposing complexity,
generating hypotheses, evaluating evidence, and synthesising a path forward.
AI can participate in every stage of this process, but only as a collaborator.
The structured thinking must come from you.

\section{Structured Problem Decomposition}
\label{sec:decomposition}

The most powerful use of AI in problem-solving is not asking it for an answer ---
it is asking it to help you structure the problem before you try to solve it.

\textbf{Decomposition prompts:}
\begin{itemize}
  \item \textit{``Break this problem into its independent sub-problems.''}
  \item \textit{``What are the three biggest uncertainties in this situation?''}
  \item \textit{``What do I need to know before I can make a good decision here?''}
  \item \textit{``Draw a cause-and-effect chain from the observed symptom to the likely root causes.''}
\end{itemize}

AI is particularly good at generating a structured first map of a problem.
You then critique the map, add what it missed, and remove what is not relevant.

\section{The AI Rubber Duck Technique}
\label{sec:rubberduck}

Rubber duck debugging --- explaining a problem out loud until the act of explanation
reveals the solution --- works dramatically better with AI than with a real rubber duck,
because AI responds.

\textbf{How to use it:}
\begin{enumerate}
  \item Describe the problem in plain language to AI as if explaining it to a
        non-expert colleague.
  \item Let AI ask clarifying questions or point out what you haven't explained.
  \item Often, the act of constructing a clear explanation surfaces the answer.
  \item If not, AI's follow-up questions tend to expose the assumption or gap you
        were not seeing.
\end{enumerate}

\begin{tipbox}
When you are stuck, write to AI as if writing to a thoughtful colleague who knows
nothing about your specific system.
The constraint of complete explanation often breaks the mental block.
\end{tipbox}

\section{Root Cause Analysis}
\label{sec:rca}

AI is useful for structuring root cause analysis (RCA) after an incident or failure.
It knows common frameworks (5 Whys, Fishbone, Fault Tree) and can apply them to
your scenario.

\textbf{RCA prompts:}
\begin{itemize}
  \item \textit{``Apply the 5 Whys to this incident: [describe]. What is the likely systemic root cause?''}
  \item \textit{``Generate a Fishbone diagram structure for this failure [describe] across these categories: People, Process, Technology, Environment.''}
  \item \textit{``What systemic prevention measures would address these root causes, beyond just fixing the immediate trigger?''}
\end{itemize}

\section{Generating and Evaluating Options}
\label{sec:options}

One of AI's strongest roles in problem solving is generating options you haven't
thought of, and then pressure-testing each option against your constraints.

\begin{enumerate}
  \item \textit{``Give me five different ways to solve [problem], ranging from minimal effort to comprehensive.''}
  \item \textit{``What are the pros, cons, and risks of each option?''}
  \item \textit{``Which option would you recommend given these constraints: [list], and why?''}
  \item \textit{``What's the failure mode of the recommended option?''}
\end{enumerate}

\begin{warnbox}[Overconfidence Risk]
AI presents all options with similar confidence regardless of how viable they are.
A bad option and a good option look identical in format.
You must apply domain knowledge to evaluate which options are actually feasible
in your context.
\end{warnbox}

\section{When AI Hurts Problem Solving}
\label{sec:ai-hurts}

AI can actively harm problem-solving when:

\begin{itemize}
  \item \textbf{You outsource thinking too early.} Jumping to AI before forming your
        own hypothesis degrades your own analytical skill over time.
  \item \textbf{You accept the first answer.} The first answer is a starting point.
        Push back, ask for alternatives, challenge the reasoning.
  \item \textbf{The problem requires tacit knowledge.} Organisation-specific, cultural,
        or relational problems cannot be solved by an AI that doesn't know your organisation.
  \item \textbf{You mistake plausibility for correctness.} AI generates fluent,
        confident-sounding text. Plausible is not the same as accurate.
\end{itemize}

Use AI in the later stages of problem framing and in option generation, not as a
replacement for your own initial analysis.
